% Generated draft: requires mathematical and source review.
\documentclass[11pt]{article}
\usepackage[margin=25mm]{geometry}
\usepackage{fontspec}
\setmainfont{Latin Modern Roman}
\usepackage{amsmath,amssymb,mathtools,unicode-math}
\setmathfont{Latin Modern Math}
\usepackage{enumitem,needspace,xurl,hyperref}
\setlength{\parindent}{0pt}
\setlength{\parskip}{5pt}
\setlength{\emergencystretch}{3em}
\newcommand{\N}{\mathbb{N}}
\newcommand{\Z}{\mathbb{Z}}
\newcommand{\Q}{\mathbb{Q}}
\newcommand{\R}{\mathbb{R}}
\newcommand{\C}{\mathbb{C}}
\newcommand{\F}{\mathbb{F}}
\newcommand{\A}{\mathbb{A}}
\newcommand{\PP}{\mathbb{P}}
\newcommand{\E}{\mathbb{E}}
\newcommand{\eps}{\varepsilon}
\newcommand{\dd}{\,\mathrm{d}}
\begin{document}
\section*{UnsolvedMath ID: 30006330}
\textbf{Source problem code:} catalogue entry
\subsection{Definitions and mathematical statement}
The supplied fragment does not determine a fully specified minimization problem or a conjectured conclusion.

To describe the missing data, let \(X\) denote an unspecified ambient space, let \(\mathcal A\) denote an unspecified class of admissible subsets of \(X\), and let \(\mathcal F\) denote an unspecified objective functional on \(\mathcal A\). If a group \(G\) acts on \(X\), a set \(E\subseteq X\) is \(G\)-invariant when
\[
 gE=E\qquad\text{for every }g\in G,
 \qquad gE:=\{g\cdot x:x\in E\}.
\]
If imposing symmetry means requiring each competitor to be invariant, the associated schematic problems would be
\[
 \inf_{E\in\mathcal A}\mathcal F(E)
 \qquad\text{and}\qquad
 \inf_{\substack{E\in\mathcal A\\ gE=E\ \text{for all }g\in G}}\mathcal F(E).
\]
These expressions illustrate one possible interpretation; the source does not specify that this is the intended meaning of imposing symmetry. In particular, invariance of each competitor and invariance of the objective under a group action are different requirements.

\textbf{Source question, with unresolved qualifications.} What effect does a discrete group of symmetries have on the intended minimization problem over sets, and does the intended analogous conclusion hold upon restricting the competitors to polygons in two dimensions?

The ambient space, functional, admissibility constraints, symmetry group and action, and meaning of discreteness are absent. The source also does not state which conclusion about minimization is sought. Consequently, neither a definite comparison nor a quantified conjecture can be recovered. The planar restriction supplies the dimension and the word ``polygons,'' but does not specify the polygon class or the conclusion to be transferred. No existence, regularity, symmetry, or polygonality assertion is made here.
\subsection{Short English statement}
The fragment asks how discrete symmetries affect minimization over sets and whether an unspecified analogous conclusion holds for planar polygons. The objective, constraints, symmetry action, polygon class, and intended conclusion are missing, so a precise conjecture requires further source clarification. The supplied open label is dated curation, not an independently verified status.
\subsection{Sources}
\begin{itemize}
\item Supplied local UnsolvedMath registry record, permanent ID \(30006330\), title \emph{Discrete Symmetries in Geometric Set Minimization}, source problem code \texttt{catalogue entry}, in the Oberwolfach Reports Open Problems collection.
\item Supplied report citation: \emph{Singularities in Discrete Systems} (2025), Oberwolfach Reports, DOI 10.4171/owr/2025/23, pp.~1163--1214. \url{https://doi.org/10.4171/owr/2025/23}. This is the supplied authoritative report reference; no exact location within that page range is given for the fragment.
\item Supplied attribution attached to the original fragment: ``Discreteness makes you think,'' Riccardo Cristoferi, joint work with Marco Bonacini and Ihsan Topaloglu. The registry gives no separately identified proposer; the relationship between this attribution and the precise question requires source review.
\item Supplied literature reference: \emph{Minimality of Polytopes}. \url{https://arxiv.org/abs/2004.03628}. Its relevance to a particular functional or theorem is not specified in the supplied record.
\item Literature assessment checked at \texttt{2026-08-22}: the supplied curation labels the question open and states that existing work proves only restricted small-mass and symmetry cases. This is supplied dated curation, not an independently verified literature assessment. The background explicitly requests independent verification of time-sensitive status and cited claims.
\item The registry labels the statement \texttt{corrected\_verified} and reports that the official Oberwolfach report was checked, the question isolated, and notation repaired. These are supplied source-verification claims; no independent source consultation or verification is asserted in this draft.
\end{itemize}
\end{document}
